\documentclass[11pt]{article}
\usepackage[margin=1in]{geometry}
\usepackage{booktabs}
\usepackage{hyperref}
\usepackage{microtype}
\usepackage{amsmath}

\title{EvidRoute: Risk-Constrained Sequential Routing over\\
Heterogeneous Evidence and Interaction Sources}
\author{Pankit Brahmkhatri\\Independent Researcher}
\date{Engineering prototype --- public benchmark experiments unrun}

\begin{document}
\maketitle

\begin{abstract}
Retrieval-augmented systems commonly decide whether to search, but operational agents may
choose among parametric knowledge, memory, lexical and dense indices, structured sources,
frozen or live web evidence, clarification, and abstention. We present EvidRoute, an
offline-first testbed that formulates this choice as bounded sequential evidence acquisition
under cost, latency, privacy, and selective-risk constraints. The current release implements
typed routes, forced-route potential outcomes, a transparent sequential policy, a CPU learned
router, exact provenance, deterministic shifts, and replayable traces. Results in this artifact
are limited to the synthetic MiniRoute smoke benchmark; external benchmark evaluation remains
future work.
\end{abstract}

\section{Introduction}
Adaptive retrieval can reduce unnecessary search, but uncertainty about when to retrieve does
not resolve which source to use or when to ask the user. Prior work studies uncertainty-aware
retrieval \cite{moskvoretskii2025adaptive,wu2025search}, information-foraging policies
\cite{qian2025scent}, and knowledge-grounded conversational agents
\cite{shi2026tauknowledge}. EvidRoute isolates heterogeneous route selection as the experimental
object.

\section{Problem}
At step $t$, a policy observes query features, remaining budget $b_t$, source-health state, and
previous evidence. It selects a feasible route $r_t$ or a terminal action. We optimize
\[
U = q_{\mathrm{supported}} + 0.35 q_{\mathrm{action}}
    - 0.15 c - 0.08\log(1 + \ell/100) - 0.45 q_{\mathrm{contradiction}},
\]
subject to a target upper bound on unsupported-or-incorrect answer risk. The coefficients are
configurable normative weights, not learned facts.

\section{System}
Routes share availability, estimate, probe, acquire, health, and close operations. Acquired
content is normalized into evidence records with source URI, snapshot, retrieval score,
timestamp, privacy class, integrity hash, and relation path. Verified answers require citations.
Clarification and abstention are first-class outcomes.

\section{Experimental protocol}
MiniRoute contains synthetic direct, lexical, semantic, structured, temporal, memory,
contradiction, ambiguity, privacy, injection, outage, and shift cases. Forced-route runs produce
potential outcomes and an oracle. The learned router predicts route utility from decision-time
features. Public evaluation is pre-specified for HotpotQA \cite{yang2018hotpotqa}, MuSiQue
\cite{trivedi2022musique}, and 2WikiMultiHopQA \cite{ho2020wiki}; those runs are not reported in
this release.

\section{Risk and shift}
The controller selects a score threshold whose one-sided Wilson upper bound is below the target.
This is an auditable prototype consistent with risk-controlling prediction principles
\cite{bates2021distribution}. Detected source-score or error-rate shift invalidates the
exchangeability status and triggers abstention or recalibration.

\section{Current evidence}
The executed MiniRoute smoke run validates software integration and artifact generation only.
It does not establish superiority, broad generalization, or production safety. Public benchmark
tables are intentionally omitted until reproducible runs exist.

\section{Limitations}
The synthetic benchmark is small and English-only. Dense retrieval is a deterministic CPU
fixture, the parametric route is a mock, and live web is disabled. Utility weights embed design
preferences. Prompt-injection detection is incomplete. Calibration guarantees do not survive
arbitrary distribution shift.

\section{Ethics}
The system exposes provenance, privacy denials, conflicts, and abstentions, but users must still
validate sources and protect data. Private $\tau$-Knowledge materials are never redistributed.
No institutional affiliation or state-of-the-art claim is asserted.

\bibliographystyle{plain}
\bibliography{references}
\end{document}
